\documentclass[a4paper,11pt]{article}

\usepackage[T1]{fontenc}
\usepackage[french]{babel}
\usepackage{amsmath}
\usepackage{amssymb}
\usepackage{geometry}

\geometry{margin=2.5cm}

\title{Formules et références}
\author{Prénom Nom}
\date{\today}

\begin{document}

\maketitle

\section{Équations en ligne}

Complétez cette section avec deux équations en ligne.

\begin{align}
    x &= x + 1
    \\
    y &= y + 1
\end{align}

\section{Équations numérotées}

Ajoutez une équation numérotée avec un label, puis citez-la dans le texte.
\begin{equation}
    a + b = c 
    \label{eq:addition} % Ce label permet ensuite de citer l'équation par son numéro dans un text avec la commande \eqref{}
\end{equation}

L'équation \eqref{eq:addition} est une addition

\section{Système}

Ajoutez un système d'équations avec l'environnement \texttt{cases}.
\[
\begin{cases}
    x = x+1
    \\
    y = y+2
\end{cases}
\]

\section{Alignement}

Ajoutez un calcul sur plusieurs lignes avec l'environnement \texttt{align}.

\begin{align}
    x &= 12 + 2 \times 5 
    \\
      &= 12 + 10 
    \\
      &= 22
\end{align}

\end{document}
